%---------------------------------------------------------------------
\section{Recovery}
\label{sec:recovery}

\subsection*{Preamble: Climbing-Specific Recovery Constraints}

Rock climbing imposes two distinct recovery clocks that must not be conflated. Forearm flexor \emph{muscle} recovers in 24--72~h via standard oxidative and glycolytic replenishment pathways. Finger flexor \emph{tendons and annular pulleys} turn over collagen on a timeline of days to weeks; A2 and A4 pulley loading during crimp grip has been modelled at 360--450~N \cite{VQLM06}. Recovery modalities must therefore be evaluated separately against each target tissue, and whether reported outcomes are surrogate (biomarkers, DOMS scores, dynamometric) or functional (climbing performance, injury incidence) must be stated explicitly throughout.

No RCT has tested any recovery modality specifically against climbing performance or pulley/tendon injury incidence as a primary outcome. Every extrapolation from general athletic populations is noted as \textbf{Bioplausible} rather than \textbf{Evidence-backed}.

%---------------------------------------------------------------------
\subsection*{Summary Table of Cited Peer-Reviewed Studies}

\begin{center}
\scriptsize
\setlength{\tabcolsep}{1pt}
\begin{longtable}{@{}>{\raggedright}p{3.0cm} >{\raggedright}p{3cm} >{\raggedright}p{2.5cm} >{\raggedright}p{2.5cm} >{\raggedright}p{2.5cm} >{\raggedright}p{1.6cm} >{\raggedright\arraybackslash}p{1.8cm}@{}}
\toprule
\textbf{Study} & \textbf{Design / n} & \textbf{Intervention} & \textbf{Primary Outcome} & \textbf{Effect Size} & \textbf{Surr./ Funct.} & \textbf{Quality} \\
\midrule
\endfirsthead
\toprule
\textbf{Study} & \textbf{Design / n} & \textbf{Intervention} & \textbf{Primary Outcome} & \textbf{Effect Size} & \textbf{Surr./ Funct.} & \textbf{Quality} \\
\midrule
\endhead
\midrule
\multicolumn{7}{r}{\footnotesize\textit{Continued on next page}}\\
\endfoot
\bottomrule
\endlastfoot

Mah et al.\ (2011). \emph{Sleep} 34(7):943--950. \cite{MMKD11}
& Prospective cohort; n=11 collegiate basketball players; no control group
& Sleep extension \textasciitilde{}+110~min/night over 5--7~wk vs.\ habitual
& Sprint time, shooting accuracy, PVT reaction time
& Sprint 16.2$\to$15.5~s ($\Delta$=0.7~s, p<0.001); shooting +9\%; PVT 310.8$\to$274.5~ms (p=0.004)
& Functional (sprint, shooting); Surrogate (PVT)
& Moderate: no control; sport-specific \\

Milewski et al.\ (2014). \emph{J Pediatr Orthop} 34(2):129--133. \cite{MSBE14}\footnote{Single-source claim --- requires independent verification.}
& Cross-sectional; n=112 adolescent athletes; 1 season follow-up
& Sleep $<$8~h vs.\ $\geq$8~h per night
& Sport injury incidence over season
& OR 1.70 (95\% CI 1.02--2.84) for injury with $<$8~h sleep
& Functional (injury incidence)
& Moderate: observational; confounders partially controlled \\

Halson (2014). \emph{Sports Med} 44(Suppl~1):13--23. \cite{HAL14}
& Narrative review; elite athletes
& Sleep duration and hygiene
& Performance impairment with $<$7~h; recovery markers
& No pooled ES; consensus: sleep deprivation $\downarrow$ reaction time, strength, glycogen resynthesis
& Surrogate and Functional
& Review; no pooled ES \\

Kiviniemi et al.\ (2007). \emph{Eur J Appl Physiol} 101(6):743--751. \cite{KHKT07}
& RCT; n=26 moderately fit males; 4~wk
& HRV-guided vs.\ fixed-schedule endurance training
& VO\textsubscript{2}peak, maximal running velocity (Loadmax)
& HRV group: $\Delta$VO\textsubscript{2}peak +3.9 vs.\ +1.4~ml\textperiodcentered{}kg\textsuperscript{-1}\textperiodcentered{}min\textsuperscript{-1} (p=0.002); Loadmax +0.9$\pm$0.2 vs.\ +0.5$\pm$0.4~km\textperiodcentered{}h\textsuperscript{-1} (p=0.048)
& Functional
& Moderate: small n; 4~wk; endurance population only \\

Dupuy et al.\ (2018). \emph{Front Physiol} 9:403. \cite{DDTE18}
& Systematic review + network meta-analysis; 99 studies; n=1,395
& Massage, CWT, CWI, compression, active recovery vs.\ passive rest
& DOMS, perceived fatigue, inflammation, performance surrogates
& Massage: SMD $-$0.72 (95\% CI $-$1.13 to $-$0.31) DOMS; SMD $-$0.47 ($-$0.76 to $-$0.18) fatigue. CWT for DOMS: SMD range $-$0.40 to $-$0.60 across contributing trials (heterogeneous protocols)
& Surrogate (DOMS, fatigue); some functional
& High: network MA; outcomes predominantly perceptual \\

Roberts et al.\ (2015). \emph{J Physiol} 593(18):4285--4301. \cite{RRME15}
& Parallel RCT; n=21 resistance-trained males; 12~wk strength training
& CWI (10°C, 10~min) vs.\ active recovery post-session
& Muscle CSA, satellite cell activity, anabolic signalling (p70S6K), isokinetic strength at 12~wk
& Type~II fibre CSA: +17\% ACT vs.\ \textasciitilde{}0\% CWI; isokinetic work: +19\% ACT only; p70S6K phosphorylation blunted with CWI (p<0.05)
& Surrogate (CSA, signalling) + Functional (strength)
& High: mechanistic + functional; lower limb --- extrapolation to forearm cautious \\

Pearcey et al.\ (2015). \emph{J Athl Train} 50(1):5--13. \cite{PBKE15}
& RCT crossover; n=8
& Foam rolling 20~min post-exercise vs.\ passive rest
& DOMS (VAS), sprint time, jump, strength-endurance at 24--72~h
& Sprint 24~h: $-$4.4\% vs.\ control; DOMS $d$=0.59--0.84; performance $d$=0.48--0.87
& Surrogate (DOMS) + Functional (sprint)
& Moderate: small n; lower extremity only \\

Molouki et al.\ (2016). \emph{J Phys Ther Sci} (forearm massage RCT)\footnote{Single-source claim --- requires independent verification.}
& RCT; n=44 healthy males
& 5~min forearm/hand massage vs.\ passive movement
& Grip endurance, grip strength (surrogate for climbing endurance)
& Endurance: 38.4$\to$46.5~s (+21\%, p<0.001); strength change NS (p=0.077)
& Surrogate/Functional (grip endurance)
& Moderate: single session; generalisation to repeated climbing loads uncertain \\

Mourot et al.\ (2015). \emph{J Sport Sci Med} 14:769--776.\footnote{Single-source claim --- requires independent verification.}
& Randomised crossover; n climbers (sport-specific)
& Active recovery climbing vs.\ walking recovery between maximal bouts
& Subsequent climbing performance (metres ascended)
& $d$=0.6 (p<0.01) favouring active climbing recovery
& Functional (climbing performance)
& Moderate: climbing-specific; small n \\

Shaw et al.\ (2017). \emph{Am J Clin Nutr} 105(1):136--143. \cite{SLRE17}
& Crossover RCT; n=8 healthy males; double-blind
& 15~g gelatin + 48~mg vitamin~C vs.\ placebo, 1~h pre-intermittent skipping exercise
& Circulating P1NP (procollagen type~I N-terminal propeptide); engineered ligament collagen content
& P1NP AUC \textasciitilde{}2$\times$ placebo ($d\approx$1.2 from published figure; 95\% CI not reported); ligament collagen content and mechanics improved
& Surrogate (biomarker P1NP; engineered tissue)
& Moderate: small n; surrogate biomarker; no tendon injury outcome; no climbing population \\

Moore et al.\ (2009). \emph{Am J Clin Nutr} 89(1):161--168. \cite{MRFE09}
& Dose-response RCT; n=6
& Whey protein 0--40~g post-resistance exercise
& Myofibrillar protein synthesis (isotope tracer MPS)
& MPS plateaus at \textasciitilde{}0.24~g/kg per meal; 20~g sufficient for \textasciitilde{}80~kg athlete
& Surrogate (isotope tracer MPS)
& High: mechanistic; isotope tracer; functional translation confirmed in longer-term RCTs \\

\end{longtable}
\setlength{\tabcolsep}{6pt}
\captionof{table}{Summary of peer-reviewed studies cited in this section. ES = effect size; CWT = contrast water therapy; CWI = cold-water immersion; DOMS = delayed-onset muscle soreness; P1NP = procollagen type~I N-terminal propeptide; MPS = myofibrillar protein synthesis.}
\end{center}

%---------------------------------------------------------------------
\subsection{Sleep: Duration, Quality, and Climbing Performance}

\begin{evidencebox}{Sleep Extension and Athletic Performance: Basketball Cohort Evidence · EVIDENCE\_BACKED · ESS 4}
\textbf{Sleep extension improves athletic performance metrics (basketball, non-climbing population).} Mah et al.\ (2011) extended sleep by +110~min/night in collegiate basketball players over 5--7~weeks and observed: sprint time 16.2$\to$15.5~s ($\Delta$=0.7~s, p<0.001); free-throw shooting +9\%; PVT mean reaction time 310.8$\to$274.5~ms (p=0.004). Outcomes are functional (sprint, shooting) and surrogate cognitive (PVT). No control group; causal inference is limited but effect magnitudes are large and directionally consistent with laboratory sleep deprivation work. \cite{MMKD11}

\textbf{Sleep $<$8~h/night associates with elevated injury risk.}\footnote{Single-source claim --- requires independent verification.} Milewski et al.\ (2014) found adolescent athletes sleeping $<$8~h faced OR~1.70 (95\% CI 1.02--2.84) for sport injury over a season vs.\ $\geq$8~h; confounders partially controlled. Functional outcome (injury incidence). \cite{MSBE14}

\textbf{Halson (2014) narrative review} concluded that elite athletes averaging $<$7~h show impaired reaction time, strength, and endurance. No pooled effect size. \cite{HAL14}
\end{evidencebox}

\begin{bioplausbox}{Sleep and Connective Tissue Collagen Synthesis: SWS--GH--IGF-1 Pathway · BIOPLAUSIBLE · ESS 2}
\textbf{Sleep and connective tissue collagen synthesis.} During slow-wave sleep (SWS, NREM~stage~3), pulsatile GH secretion accounts for the majority of daily GH output (Van Cauter et al., 2000). GH stimulates hepatic IGF-1, which upregulates collagen gene expression in tenocytes and fibroblasts. Chronic sleep restriction reduces SWS and therefore GH pulse amplitude. \emph{Direct RCT evidence linking sleep duration or SWS quality to tendon or pulley collagen synthesis, repair timeline, or injury incidence in humans is absent.} Mechanism: SWS $\to$ GH $\to$ IGF-1 $\to$ collagen type~I gene upregulation in fibroblasts/tenocytes; gap: no human tendon biopsy or ultrasound-assessed tendon morphology under controlled sleep manipulation.

Saner et al.\ (2020; \textbf{TO\_VERIFY: unresolved source DOI in CrossRef lookup}) reported acute sleep deprivation reduced muscle protein synthesis by 18\% in humans (surrogate metabolic outcome, isotope tracer); the extent to which this applies to the slower tendon collagen synthesis pathway is unknown.
\end{bioplausbox}

\begin{folklorebox}{Sleep Hygiene Practices for Athletes · FOLKLORE\_VERIFIED · ESS 1}
\textbf{Sleep hygiene practices} (dark, cool room; consistent wake time $\pm$30~min; minimize screen-based blue light $\geq$60~min pre-sleep; avoid caffeine $>$6~h pre-sleep; 20--30~min nap to supplement poor nights) are standard sport-science coaching consensus without climbing-specific RCT validation. Sources: Halson (2014) review; US Olympic Committee sleep guidelines; circadian biology literature. Blue-light avoidance is underpinned by melatonin suppression mechanism (\textbf{Bioplausible}).
\end{folklorebox}

\textbf{Practical targets:} $\geq$8~h/night; aim for 9~h time-in-bed during high-volume training blocks. \textbf{Evidence-backed} for general athletic performance (Mah 2011); \textbf{Bioplausible} for connective tissue benefit. No climbing-specific sleep RCT exists as of 2024 --- a critical evidence gap.

\textbf{Critical note for connective tissue:} Any practice reducing SWS quality --- alcohol, late caffeine, chronic high training stress --- may impair connective tissue repair even if total sleep duration is nominally adequate. This pathway deserves direct investigation in climbers.

%---------------------------------------------------------------------
\subsection{HRV Monitoring as Readiness Proxy}

\begin{evidencebox}{HRV-Guided Training: Endurance Population RCT Evidence · EVIDENCE\_BACKED · ESS 4}
\textbf{HRV definition and metric.} HRV quantifies beat-to-beat R--R interval variation derived from ECG or photoplethysmography. RMSSD (root mean square of successive differences) is the primary parasympathetic index recommended for daily athlete monitoring; ln(RMSSD) improves normality and is used in HRV4Training (Plews et al., 2013, \cite{WHHE13}, case series n=4 elite triathletes --- low evidence tier for decision rules, but establishes metric rationale).

\textbf{HRV-guided vs.\ fixed training (only RCT evidence).} Kiviniemi et al.\ (2007) randomised 26 recreational runners to HRV-guided training (intensity adjusted daily based on morning HRV) or fixed periodised schedule over 4~weeks. HRV-guided group: $\Delta$VO\textsubscript{2}peak +3.9 vs.\ +1.4~ml\textperiodcentered{}kg\textsuperscript{-1}\textperiodcentered{}min\textsuperscript{-1} (p=0.002); Loadmax +0.9$\pm$0.2 vs.\ +0.5$\pm$0.4~km\textperiodcentered{}h\textsuperscript{-1} (p=0.048). Functional outcome; endurance population; 4-week duration; small n. \cite{KHKT07}

\textbf{Device measurement validity.} HRV4Training smartphone PPG agrees with ECG-derived RMSSD within acceptable limits (Altini \& Kinnunen, 2021, \emph{Sensors} 21(12):4124, \cite{MS21}). Measurement validity does not equal readiness-score validity; proprietary algorithms (Whoop, Garmin Body Battery) have not been validated against performance or injury outcomes in climbing populations.
\end{evidencebox}

\begin{bioplausbox}{HRV as Tendon/Pulley Readiness Proxy: Mechanistically Unjustified · BIOPLAUSIBLE · ESS 2}
\textbf{HRV as tendon/pulley readiness proxy: mechanistically unjustified.} HRV reflects autonomic nervous system balance, not local connective tissue repair status. Tendon and pulley tissue lacks the direct autonomic innervation that would create an HRV signal. A climber may register ``green'' HRV while an A2 pulley remains mechanically sensitised from the prior session. \emph{No study has compared HRV metrics to tendon mechanobiological markers, ultrasound-assessed pulley/tendon morphology, or finger injury risk in climbers.} Use HRV to detect systemic overreaching and neuromuscular/cardiovascular under-recovery, not to clear finger tissues for load.
\end{bioplausbox}

\begin{toverifybox}{HRV ``Ready to Train'' Thresholds for Climbing · TO\_VERIFY · ESS 1}
\textbf{TO\_VERIFY --- ``Ready to train'' thresholds:} (e.g., daily RMSSD within $\pm$0.5~SD of 7-day rolling mean = proceed; $>$1~SD drop = reduce intensity) derive from empirical conventions developed from Ithlete and HRV4Training user data and endurance coaching practice (Plews et al., Buchheit HRV+ framework). These are not derived from RCTs with performance or injury endpoints as primary outcomes in strength/isometric sports. \emph{No universal, sport-validated cutoff exists.} Requires $\geq$3 independent climbing-specific sources for FOLKLORE\_VERIFIED status.
\end{toverifybox}

\textbf{Practical use in climbers:} Track morning RMSSD against a 7--10~day rolling baseline. Use suppressions as a secondary systemic fatigue signal alongside: (1) local finger pain/morning stiffness, (2) skin status, (3) hang warm-up quality, (4) session RPE. A normal HRV does not imply pulley readiness; finger symptoms override HRV.

%---------------------------------------------------------------------
\subsection{Contrast Therapy and Cold-Water Immersion}

\subsubsection*{Mechanism}

\begin{bioplausbox}{Contrast Water Therapy: Vascular Pump Mechanism · BIOPLAUSIBLE · ESS 2}
Alternating cold (8--15°C) and warm (36--42°C) water immersion is theorised to create cyclical vasoconstriction--vasodilation, acting as a ``vascular pump'' to clear metabolic waste products, reduce oedema, and enhance local perfusion (Bleakley \& Davison, 2010, \emph{BJSM} 44(3):179--187, \cite{BD09}). Direct controlled evidence for this specific pump mechanism in the forearm is absent; the rationale is mechanistically plausible.
\end{bioplausbox}

\subsubsection*{Evidence for Muscle Soreness Recovery}

\begin{evidencebox}{Contrast Water Therapy Reduces DOMS (Dupuy 2018 Network Meta-Analysis) · EVIDENCE\_BACKED · ESS 4}
Dupuy et al.\ (2018) network meta-analysis (99 studies, n=1,395): CWT reduces DOMS vs.\ passive rest; SMD values from contributing trials range approximately $-$0.40 to $-$0.60 (heterogeneous protocols across sources; exact pooled CWT-specific estimate was not consistently reported with CIs across all contributing sources). The effect on \emph{performance} outcomes (strength, power, sprint) was less consistent and smaller. Outcomes are primarily surrogate (DOMS, perceived fatigue). \cite{DDTE18}
\end{evidencebox}

\subsubsection*{Critical Concern: CWI Blunts Strength Adaptation}

\begin{evidencebox}{Cold-Water Immersion Blunts Strength Adaptation (Roberts 2015) · EVIDENCE\_BACKED · ESS 4}
Roberts et al.\ (2015) RCT (n=21, 12-week lower-body strength training programme) demonstrated that CWI (10°C, 10~min) post-session significantly attenuated: type~II fibre CSA (+0\% vs.\ +17\% active recovery; p<0.05); satellite cell proliferation; p70S6K phosphorylation (blunted at 24 and 48~h); isokinetic work (+0\% vs.\ +19\% active recovery). Outcomes: surrogate (CSA, anabolic signalling) and functional (strength). Lower-limb population --- extrapolation to forearm flexors is cautious but mechanistically justified given shared mTOR/satellite cell pathway. \textbf{Routine post-hangboard or post-strength CWI likely blunts intended forearm flexor adaptation.} \cite{RRME15}
\end{evidencebox}

\subsubsection*{Connective Tissue Recovery}

\begin{bioplausbox}{CWT/CWI and Tendon Collagen Synthesis: Theoretical Anti-Inflammatory Concern · BIOPLAUSIBLE · ESS 2}
No published RCT has examined CWT or CWI effects on tendon or pulley collagen synthesis or repair. The anti-inflammatory effect of CWI could theoretically suppress the prostaglandin-mediated phase of tendon remodelling; if this phase is necessary for tendon adaptation, CWI might impair rather than accelerate repair. Gap: no human tendon biopsy or ultrasound data under CWI conditions post-loading.
\end{bioplausbox}

\textbf{Practical protocol (if used):}
\begin{itemize}[noitemsep]
  \item \textbf{Avoid routine CWI immediately post-hangboard or strength session} to preserve anabolic signalling. \textbf{Evidence-backed} (Roberts 2015).
  \item CWT is appropriate after technique or high-volume/pump days where hypertrophic adaptation is not the goal: cold phase 10--12°C, 1--2~min; hot phase 38--40°C, 1--2~min; 3--4 cycles; $\leq$2$\times$/week. \textbf{Evidence-backed} for DOMS reduction (Dupuy 2018); \textbf{Bioplausible} for forearm-specific application.
\end{itemize}

%---------------------------------------------------------------------
\subsection{Soft-Tissue Work: Massage and Foam Rolling for Forearm Recovery}

\subsubsection*{General Evidence}

\begin{evidencebox}{Massage and Foam Rolling for Muscle Soreness Recovery (Dupuy 2018; Pearcey 2015) · EVIDENCE\_BACKED · ESS 4}
\textbf{Massage.} Dupuy et al.\ (2018) network meta-analysis: massage showed the largest effect among all modalities for muscle soreness reduction, SMD~$=$~$-$0.72 (95\% CI $-$1.13 to $-$0.31), and for perceived fatigue reduction, SMD~$=$~$-$0.47 (95\% CI $-$0.76 to $-$0.18). Outcomes are predominantly surrogate/perceptual; direct performance recovery effects were smaller and less consistent. \cite{DDTE18}

\textbf{Foam rolling.} Pearcey et al.\ (2015) crossover RCT (n=8): 20~min post-exercise foam rolling reduced DOMS ($d$=0.59--0.84) and attenuated performance decrements at 24~h (sprint $-$4.4\% vs.\ control; $d$=0.48--0.87). Lower extremity; surrogate + functional outcomes. \cite{PBKE15}
\end{evidencebox}

\subsubsection*{Forearm-Specific Evidence}

\begin{evidencebox}{Forearm Massage Improves Grip Endurance (Molouki 2016) · EVIDENCE\_BACKED · ESS 4}
\textbf{Forearm massage (single-session RCT).}\footnote{Single-source claim --- requires independent verification.} Molouki et al.\ (2016) RCT (n=44 healthy males): single 5-min forearm/hand massage increased grip endurance from 38.4 to 46.5~s (+21\%, p<0.001); grip strength change was non-significant (p=0.077). Outcome is a functional surrogate for climbing endurance. Limitations: single session, no repeated-loading context.
\end{evidencebox}

\begin{bioplausbox}{Forearm Soft-Tissue Work for Climbing Recovery: Mechanism Without Climbing RCT · BIOPLAUSIBLE · ESS 2}
No RCT has evaluated forearm massage or foam rolling specifically after climbing-equivalent isometric or dynamic loading. Proposed mechanisms --- increased local blood flow, reduced myofascial stiffness, parasympathetic activation, interoceptive modulation --- are plausible but unverified for the forearm flexor compartment in climbers. Soft-tissue work does not directly alter tendon collagen synthesis rates or structural integrity.
\end{bioplausbox}

\begin{folklorebox}{Climbing-Specific Forearm Recovery Tools · FOLKLORE\_VERIFIED · ESS 1}
\textbf{Climbing-specific forearm tools:} forearm rolling with lacrosse ball or dowel, finger flexor release with thumb pressure, massage gun to forearm bellies, ``voodoo floss'' occlusion-reperfusion. Sources: Lattice Training programming, Dave MacLeod coaching resources, Eva~L\'opez training recommendations, Esther Smith / The Climb Clinic physiotherapy community. No controlled trials exist for these specific applications.
\end{folklorebox}

\textbf{Contraindication:} Aggressive deep-tissue massage or foam rolling directly over an acutely loaded A2 pulley is not supported by any evidence and is potentially harmful. The inflammatory phase of tendon/pulley remodelling serves biological functions; disruption via aggressive mechanical pressure may impair the repair cascade. \textbf{Bioplausible concern.}

\textbf{Practical dose:} 5--15~min forearm soft-tissue work (effleurage + kneading) immediately post-session for perceived recovery and DOMS modulation; $\geq$60~s per forearm compartment. \textbf{Evidence-backed} for general DOMS reduction; \textbf{Bioplausible} for forearm-specific application.

%---------------------------------------------------------------------
\subsection{Active Recovery Sessions}

\subsubsection*{Evidence Base}

\begin{evidencebox}{Active Recovery: Lactate Clearance and Climbing-Specific Performance · EVIDENCE\_BACKED · ESS 5}
Low-intensity aerobic activity accelerates blood lactate clearance vs.\ passive rest: Cochrane (2004, \emph{J Sports Sci} 22(6):501--516) reported \textasciitilde{}30--40\% faster clearance at matched time points, via increased cardiac output sustaining oxidative phosphorylation. Surrogate outcome (lactate); mechanistically translatable to forearm but direct forearm evidence is limited.

Mota et al.\ (2017) crossover (n=14 swimmers):\footnote{Single-source claim --- requires independent verification.} 10~min active recovery at \textasciitilde{}69\% trial velocity increased lactate removal to 4.30$\pm$1.74 vs.\ 1.76$\pm$1.70~mmol\textperiodcentered{}L\textsuperscript{-1} (p<0.001) vs.\ passive rest. Surrogate outcome. \cite{MDOE17}

\textbf{Climbing-specific active recovery (single RCT).}\footnote{Single-source claim --- requires independent verification.} Mourot et al.\ (2015, \emph{J Sport Sci Med} 14:769--776) found that active recovery using climbing-specific movement between maximal bouts improved subsequent climbing performance vs.\ walking recovery ($d$=0.6, p<0.01). Functional outcome; sport-specific. This is the only climbing-specific controlled study for any recovery modality.
\end{evidencebox}

\begin{bioplausbox}{Active Recovery: Forearm Perfusion Mechanism and Conservative Load Threshold · BIOPLAUSIBLE · ESS 2}
\textbf{Forearm blood flow.} Isometric finger flexor contractions at $>$20--30\% MVC cause partial vascular occlusion; low-intensity dynamic forearm and wrist movement restores perfusion and accelerates metabolite clearance. \textbf{Lactate clearance $\neq$ connective tissue recovery:} active recovery addresses muscle readiness on a same-day or next-day timeline, not tendon collagen remodelling (days--weeks).

\textbf{The threshold between active recovery and additional connective tissue stress is not empirically defined for any climbing grade or grip position.} A conservative operational rule is: open-hand grip only, $<$40\% maximum voluntary contraction, RPE~$\leq$4/10, no small holds, no sustained isometric pump.
\end{bioplausbox}

\begin{toverifybox}{Active Recovery Formats for Climbers: Jug Laps, ARC, Antagonist Circuits · TO\_VERIFY · ESS 1}
Specific active recovery formats for climbers --- ``easy jug laps,'' ARC-style traversing, antagonist training circuits, yoga/mobility --- are codified by Anderson \& Anderson (\emph{The Rock Climber's Training Manual}, 2014) and practitioner consensus. The metabolic rationale is plausible; controlled climbing data for these specific formats is absent. Requires $\geq$3 independent named sources beyond Anderson \& Anderson for FOLKLORE\_VERIFIED status.
\end{toverifybox}

\textbf{Practical dose:} 20--30~min easy aerobic work (cycling, light jug traversing at RPE 3--4/10) on the day of or day after hard finger sessions. Antagonist and scapular stability work at low load. \textbf{Evidence-backed} for lactate clearance (extrapolated from other sports); \textbf{Bioplausible} for climbing-specific application.

%---------------------------------------------------------------------
\subsection{Nutrition Timing: Protein Dose and Timing for Connective Tissue}

\subsubsection*{Post-Exercise Protein for Muscle}

\begin{evidencebox}{Post-Exercise Protein Dose for Muscle MPS (Moore 2009; Morton 2018) · EVIDENCE\_BACKED · ESS 4}
Moore et al.\ (2009) dose-response RCT (n=6): myofibrillar MPS (isotope tracer) plateaus at \textasciitilde{}0.24~g/kg body mass per meal (approximately 20~g for an 80-kg athlete) using whey. Surrogate outcome; functional translation to hypertrophy confirmed in longer-term RCTs. \cite{MRFE09}

Morton et al.\ (2018, \emph{BJSM} 52(6):376--384, \cite{MMME17}) meta-analysis: total daily protein intake (1.6--2.2~g\textperiodcentered{}kg\textsuperscript{-1}\textperiodcentered{}day\textsuperscript{-1}) is more important than precise post-exercise timing for muscle hypertrophy. Practical target: 0.3--0.4~g/kg per meal, 3--5 meals/day.

\textbf{This evidence base is almost entirely muscle-focused. Tendon protein synthesis rates are substantially lower and driven by different stimuli (mechanical loading + amino acid substrate availability). Do not extrapolate muscle protein timing rules directly to pulley or tendon recovery.}
\end{evidencebox}

\subsubsection*{Connective Tissue-Specific: The Gelatin + Vitamin~C Protocol}

\begin{evidencebox}{Pre-Exercise Gelatin + Vitamin C: P1NP Surrogate and Ligament Mechanics (Shaw 2017) · EVIDENCE\_BACKED · ESS 4}
Shaw et al.\ (2017) crossover RCT (n=8 healthy males, double-blind): 15~g gelatin (hydrolysed collagen) + 48~mg vitamin~C vs.\ placebo ingested 1~h before 2~$\times$~6~min intermittent rope-skipping (designed to mechanically load tendons). Circulating P1NP (procollagen type~I N-terminal propeptide, marker of collagen~I synthesis) was approximately doubled in the gelatin condition vs.\ placebo at 1~h post-exercise ($d\approx$1.2 from published figure; 95\% CI not reported in manuscript). Engineered ligaments treated with post-supplement serum showed increased collagen content and mechanical stiffness. Outcomes: surrogate biomarker (P1NP) and \emph{ex vivo} engineered tissue --- no direct tendon structural outcome, no injury prevention data, no climbing population. \cite{SLRE17}

\textbf{Mechanistic rationale (Bioplausible for application):} Vitamin~C cofactors prolyl hydroxylase and lysyl hydroxylase, critical for collagen triple-helix cross-linking. Gelatin provides hydroxyproline-rich peptides (glycine, proline, hydroxyproline) that peak in circulation \textasciitilde{}60~min post-ingestion. Mechanical loading of tendons during the amino acid availability window activates mechanosensitive tenocytes, upregulating collagen gene expression. The 1-h pre-exercise timing is physiologically motivated for this cascade. Applying this to hangboard or crimp-loading sessions is mechanistically appropriate but climbing population data are absent.
\end{evidencebox}

\subsubsection*{Carbohydrate and Energy Availability}

\begin{evidencebox}{RED-S: Low Energy Availability and Injury Risk (Mountjoy 2018) · EVIDENCE\_BACKED · ESS 3}
Mountjoy et al.\ (2018, \emph{BJSM} 52(11):687--697, \cite{MSBE18}) RED-S consensus: low energy availability ($<$30~kcal/kg fat-free mass/day) is associated with increased stress fracture and tendon injury rates in endurance and aesthetic sport athletes. Mechanisms include suppressed IGF-1, sex hormones, and collagen synthesis (mechanistic evidence largely from animal models; clinical associations predominantly observational). Functional outcome (injury incidence); partially confounded.

Climbers managing weight for competition or body-mass targets are at elevated RED-S risk. \textbf{Chronic underfueling during high-load training blocks likely compromises pulley and tendon repair.} No climbing-specific RCT exists.
\end{evidencebox}

\begin{bioplausbox}{Post-Exercise Carbohydrate Co-Ingestion: Insulin-Mediated Amino Acid Uptake · BIOPLAUSIBLE · ESS 2}
Post-exercise carbohydrate co-ingestion with protein: insulin stimulates amino acid uptake in connective tissue; co-ingestion may have additive effects on tendon collagen synthesis, though this has not been directly tested in human tendon studies.
\end{bioplausbox}

%---------------------------------------------------------------------
\subsection{Practical Recovery Protocol}

For a climber training 4--5 days per week with hangboard, power-endurance, and technique sessions.

\begin{center}
\small
\begin{longtable}{>{\raggedright}p{3.2cm} >{\raggedright}p{3.5cm} >{\raggedright}p{3.0cm} >{\raggedright\arraybackslash}p{5.5cm}}
\toprule
\textbf{Timing} & \textbf{Action} & \textbf{Dose} & \textbf{Classification} \\
\midrule
\endfirsthead
\toprule
\textbf{Timing} & \textbf{Action} & \textbf{Dose} & \textbf{Classification} \\
\midrule
\endhead
\midrule
\multicolumn{4}{r}{\footnotesize\textit{Continued on next page}}\\
\endfoot
\bottomrule
\endlastfoot

Daily (every night) & Sleep window & $\geq$8~h; target 9~h during high-volume blocks; consistent wake time $\pm$30~min; dark, cool room & \textbf{Evidence-backed} for performance (Mah 2011; Milewski 2014); \textbf{Bioplausible} for connective tissue (SWS--GH--collagen axis) \\

Morning (upon waking) & HRV (RMSSD via validated device: Oura, HRV4Training) & 3--5~min seated; compare to 7--10~day rolling mean & \textbf{Evidence-backed} for neuromuscular/CV readiness (Kiviniemi 2007); \textbf{Bioplausible} for overall load management; \textbf{not} a tendon readiness proxy \\

60~min pre-session (hangboard, limit bouldering, or finger rehab) & 15~g hydrolysed collagen (gelatin) + 48--50~mg vitamin~C in water or OJ & Once per tendon-loading session & \textbf{Evidence-backed} for collagen synthesis surrogate (Shaw 2017; P1NP biomarker); \textbf{Bioplausible} for tendon/pulley functional outcome \\

Within 0--2~h post-session & 0.3~g/kg high-quality protein (whey or mixed) + 40--80~g carbohydrate if next session $<$24~h & Single meal or shake & \textbf{Evidence-backed} for muscle MPS (Moore 2009; Morton 2018) \\

Post-session (technique / power-endurance / high-volume day only) & Contrast water therapy: 3--4 cycles of cold 10--12°C / warm 38--40°C, 1--2~min each & \textasciitilde{}12--16~min total; $\leq$2$\times$/week & \textbf{Evidence-backed} for DOMS reduction (Dupuy 2018); \textbf{avoid after hangboard/strength days} (Roberts 2015) \\

Post-session (all days) & 5--15~min forearm soft-tissue work (effleurage, kneading, or foam roller/lacrosse ball on flexor bellies) & $\geq$60~s per compartment; avoid pressure directly on finger pulleys & \textbf{Evidence-backed} for general DOMS and fatigue (Dupuy 2018; Molouki 2016 for forearm grip endurance); \textbf{Bioplausible} for forearm-specific application \\

Post-hangboard / strength day & 10--20~min active recovery: light jug traversing (open hand, RPE 3--4/10), wrist mobility, light cycling. \textbf{Avoid CWI.} & Low intensity, no crimps & \textbf{Evidence-backed} for lactate clearance (extrapolated); \textbf{avoid CWI} (Roberts 2015 --- blunts hypertrophy) \\

Recovery days (1--2$\times$/week) & 20--30~min easy aerobic (cycling, walking) + scapular antagonist + mobility & RPE $\leq$4/10 & \textbf{Evidence-backed} for lactate clearance; \textbf{Bioplausible} for forearm recovery \\

Daily (all training days) & Total protein intake & 1.8--2.2~g\textperiodcentered{}kg\textsuperscript{-1}\textperiodcentered{}day\textsuperscript{-1} distributed over 4--5 meals & \textbf{Evidence-backed} for muscle adaptation (Morton 2018) \\

Daily & Energy availability & $>$40~kcal/kg fat-free mass/day; avoid chronic underfueling & \textbf{Evidence-backed} for RED-S injury associations (Mountjoy 2018) \\

\end{longtable}
\captionof{table}{Practical recovery protocol for a climber training 4--5 days per week. All recommendations filtered for finger/forearm/connective tissue relevance.}
\end{center}

\textbf{HRV decision rules (session type):}
\begin{itemize}[noitemsep]
  \item RMSSD $\geq$95\% of 7-day mean + subjectively fresh + no finger soreness: proceed with planned hangboard session.
  \item RMSSD 90--95\% of mean \textbf{or} finger morning stiffness present: reduce hangboard volume 30--50\% or substitute technique day.
  \item RMSSD $<$90\% of mean for $\geq$2 consecutive days: postpone high-intensity loading; active recovery or rest.
  \item Note: these thresholds are \textbf{Folklore} (endurance coaching consensus); not RCT-validated for isometric strength sports.
\end{itemize}

%---------------------------------------------------------------------
\subsection{What the Literature Cannot Yet Answer}

\begin{enumerate}[noitemsep]
  \item \textbf{Sleep $\times$ finger performance:} Does sleep restriction or extension directly alter finger force output, hang time-to-failure, or A2/A4 pulley repair timeline in trained climbers? The SWS--GH--IGF-1--collagen axis is biologically compelling but no study has combined polysomnography with tendon biomarkers (P1NP, collagen fractional synthetic rate) or ultrasound-assessed pulley/tendon morphology under controlled sleep manipulation in an athletic population.

  \item \textbf{HRV and finger injury prediction:} Does HRV-guided load modification reduce overuse finger injury incidence in climbers? All existing HRV-guided training RCTs use endurance populations and VO\textsubscript{2}max as the primary outcome. A multi-week RCT with A2/A4 pulley injury incidence as the endpoint would directly test whether autonomic recovery proxies benefit the sport's most prevalent injury.

  \item \textbf{CWT/CWI and tendon/pulley adaptation:} Does post-session cooling alter tendon or pulley collagen synthesis, structural stiffness, or injury risk in climbers? The muscle hypertrophy blunting effect (Roberts 2015) raises the analogous question for connective tissue where anabolic remodelling timelines are even longer.

  \item \textbf{Collagen supplementation --- functional endpoint:} Does pre-exercise gelatin + vitamin~C reduce time-loss finger tendon injuries or accelerate return-to-load after pulley injury over a full climbing season? Shaw et al.\ (2017) established an acute P1NP biomarker response; studies with structural (ultrasound, MRI) and functional (injury rate, grip force) outcomes over 8--12 repeated-dose weeks are needed.

  \item \textbf{Active recovery threshold for connective tissue:} At what intensity, grip type, and volume does ``easy'' climbing transition from promoting to impeding pulley/tendon recovery? The occlusion threshold for forearm vasculature, contraction duty cycle, and load--repair balance are unexplored experimentally in connective tissue; defining this threshold would allow evidence-based grading of recovery climbing sessions.

  \item \textbf{RED-S and finger tendons in climbers:} Does chronic low energy availability specifically impair finger flexor tendon repair in climbers, and does carbohydrate--protein co-ingestion mitigate this? Given the documented prevalence of weight-management practices in competitive climbing, the intersection of RED-S and connective tissue health is a clinically urgent, empirically unaddressed gap.
\end{enumerate}
